\section{Discussion}
\label{sec:discussion}

\subsection{Relation to existing methods}

The translation table (\cref{tab:translation}) and the RCTD worked
example (\cref{sec:rctd-worked}) connect each major ST deconvolution
method to the same coarsening framework through its parametric or
procedural choices, with RCTD derived in full and the others read off
the same C1--C2--C3 ledger: Cell2location replaces the simplex constraint
with a hierarchical prior \citep{kleshchevnikov2022cell2location};
RCTD takes the unregularized Poisson form \citep{cable2021robust};
Tangram constructs singleton candidate sets via cell-level
alignment \citep{biancalani2021deep}; CARD adds a
spatial-smoothness regularizer \citep{ma2022spatially}; SPOTlight
\citep{elosua2021spotlight}, Stereoscope
\citep{andersson2020stereoscope}, SpatialDecon
\citep{danaher2022advances}, DestVI \citep{lopez2022destvi}, and
STdeconvolve \citep{miller2022reference} are further parametric
variants. These methods perform differently on different tissues
\citep{li2023comprehensive} because each method's prior or
regularization implicitly assumes a particular candidate-set
structure; the framework-level diagnostic (rank condition +
marker-gene bias bound) lets practitioners predict which method is
appropriate before running the full pipeline.

The same signature-times-proportions algebra predates the spatial
setting in bulk RNA-seq deconvolution, where reference-based methods such
as MuSiC \citep{wang2019music}, CIBERSORTx \citep{newman2019digital}, and
the Bayesian BayesPrism \citep{chu2022bayesprism} estimate cell-type
proportions and profiles from mixed-tissue measurements. Each bulk
sample is an instance of \cref{eq:trans-spot} with a single mixture, and
a cohort of samples plays the role the spot ensemble plays here, so the
rank condition and cell-total identity carry over unchanged; what the
spatial setting adds is that the spots arrive pre-supplied with distinct
compositions from tissue architecture and that identifiability can be
restored through single-cell-resolution platforms, neither of which is
available to bulk deconvolution.

\subsection{What the framework does \emph{not} address}

\begin{itemize}
\item \textbf{Joint inference of composition and expression.} The
  framework analyzes deconvolution conditional on a cell-type
  reference. Joint inference is a richer problem
  \citep{zhang2023spatially}.

\item \textbf{Spatial autocorrelation.} The framework treats spots as
  exchangeable. CARD-style spatial-smoothness priors are an
  orthogonal addition we do not analyze.

\item \textbf{Deep-learning identifiability.} Tangram-style
  end-to-end learning may achieve identifiability for reasons not
  captured by the rank condition (e.g., implicit regularization in
  neural network optimization). Whether such methods deliver bias
  bounds comparable to \cref{thm:bias-marker} is open.

\item \textbf{Single-cell-resolution noise.} We treat MERFISH/seqFISH+
  measurements as ground-truth singletons. In practice they have
  their own noise; the framework would benefit from incorporating it
  via the C2 relaxation results of \cite{towell2026mdrelax}.

\item \textbf{External-reference dependence.} Our real-Visium
  validation uses the dataset's own CellRanger clustering as the
  cell-type reference rather than an independent scRNA-seq atlas.
  This is appropriate for testing the framework's
  conditions-on-inputs predictions (which hold regardless of
  reference choice), but a fully external reference would strengthen
  empirical claims about specific cell-type-mean recovery. The
  journal-format companion to this work pursues this with the Allen
  Brain Atlas.
\end{itemize}

\subsection{Broader implications}

The framework's main contribution is shifting the ST deconvolution
literature from method-specific to assumption-specific reasoning. A
practitioner asking ``which deconvolution method should I use?'' can
now answer by checking which assumptions (rank condition, marker-gene
adequacy, single-cell-resolution availability) hold for their data
and tissue. This is the same shift previously achieved for scRNA-seq
zero inflation in \cite{towell2026scrnacoarsening}.

\subsection{A framework across domains}

ST deconvolution is one application of a portable apparatus. The
same masked-data conditions ground a series of companion papers in
distinct domains: scRNA-seq zero inflation as a binary
candidate-set instance \citep{towell2026scrnacoarsening};
local-differential-privacy mechanisms as randomized candidate-set
generators \citep{towell2026dpcoarsening}; programmatic weak
supervision with heterogeneous labeling functions as
candidate-set unions over label space
\citep{towell2026weaksupcoarsening}; and EHR phenotyping under
code-set heterogeneity as candidate sets of clinical phenotypes
\citep{towell2026phenotypecoarsening}. Each application
contributes distinct identifiability conditions and bias bounds,
but the underlying C1, C2, C3 conditions and rank-style
identifiability criteria are shared. The portability is
itself the principal contribution of the series; this paper
demonstrates one instance.
